% !TeX root = ../main.tex
\section{How results are measured}
\label{sec:eval}

Numbers are only comparable when they use the same metric, tolerance and test
set. This section defines every metric used later and lists which of them each
of our benchmarks reports.

\subsection{Precision, recall and F1 for notes}

Compare a list of predicted notes with a list of reference (correct) notes. A
prediction \term{matches} a reference note when the pitches are equal (within
half a semitone) and the onsets are within a tolerance $\tau$. Each reference
note can be matched at most once (a one-to-one matching). Then
\[
  \text{precision} = \frac{\text{matched}}{\text{predicted}},\qquad
  \text{recall} = \frac{\text{matched}}{\text{reference}},\qquad
  F_1 = \frac{2\cdot\text{precision}\cdot\text{recall}}{\text{precision}+\text{recall}}.
\]
Precision drops when the system invents notes; recall drops when it misses
notes. F1 is high only when both are high.

\begin{worked}[: F1]
A bar has 10 reference notes. The system outputs 12 notes, 8 of which match.
Precision $= 8/12 = 0.67$, recall $= 8/10 = 0.80$,
$F_1 = 2\cdot0.67\cdot0.80/(0.67+0.80) = 0.73$.
\end{worked}

Variants used in this document:
\begin{description}[style=nextline,leftmargin=1.2cm]
\item[F50 (onset F1)] $\tau = 50\ms$, the \code{mir_eval}
  convention~\cite{raffel2014mireval,hawthorne2018onsets}. Computed on the raw
  notes of a model, before any rhythm quantisation. This is what papers report
  and what our Python stage benchmarks use.
\item[Onset-offset F1] Also requires the predicted offset to be within
  $\max(50\ms, 0.2\,d)$ of the reference offset ($d$: the reference duration).
  Measures note lengths.
\item[Final-tab F1 (ours)] $\tau = 80\ms$, computed by the Rust benchmarks on the
  notes of the finished tab, \emph{after} the notes were snapped to the rhythm
  grid. Quantisation moves onsets by up to half a grid step, hence the wider
  tolerance. Reported as the \emph{median} over items.
\item[Tab F1] A note match that is also on the same string (and hence fret) as
  the reference. Measures the whole pipeline.
\item[String accuracy] Given the \emph{correct} notes, the share placed on the
  reference string. Measures the fingering stage alone.
\item[Octave errors] The share of reference notes that were found only one or two
  octaves away. Diagnoses distortion problems.
\end{description}

\subsection{Rhythm metrics}

Beat trackers are scored with a beat \term{F-measure} (a predicted beat is
correct within 70\ms{} of a reference beat) and continuity metrics such as
\term{CMLt}, which counts only runs of consecutive correct beats, so a tracker
that keeps losing the beat scores low even if many beats are individually
right~\cite{foscarin2024beatthis}. Downbeat (first beat of each bar)
F-measure is computed the same way. Our final-tab benchmarks also count songs
whose tempo is right and whose bar lines fall on the reference bar lines.

\subsection{Separation metrics}

Separators are usually scored by the \term{signal-to-distortion ratio} (SDR),
in decibels: the energy of the reference stem divided by the energy of the
error, so higher is better and 3\dB{} more means half the error
energy~\cite{vincent2006bss}. SDR rewards a stem that \emph{sounds} close to the
reference; it does not say whether the stem is easy to transcribe. We therefore
score separators by the F50 of a transcription of their stem.

\subsection{Which number is which}

Table~\ref{tab:which-number} lists where each kind of number in this document
comes from.

\begin{table}[h]
\centering
\caption{Metrics produced by each benchmark in the repository.}
\label{tab:which-number}
\small
\begin{tabularx}{\textwidth}{@{}>{\raggedright}p{4.6cm}>{\raggedright}p{3.6cm}X@{}}
\toprule
Benchmark (code) & Metric & Notes \\
\midrule
\code{scripts/train/bench_notes.py}, \code{bench_separation.py}, \code{bench_tab.py}, \code{bench_beats.py} & F50, onset-offset F1, Tab F1, beat F / CMLt & Python, \code{mir_eval}; raw model notes; mean per item; up to 12 items per group \\
\code{bench_moisesdb.py} & F50 against notes of the isolated guitars & No human labels: ranks separators, not note models (Section~\ref{sec:bench-moises}) \\
\code{eval_guitarset} and other ignored tests in \file{transcribe/eval.rs} & Final-tab F1 (80\ms), string/fret agreement, tempo and bar lines & Rust, the real app pipeline end to end; median per item \\
\code{eval_fingering}, \code{examples/fingering.rs} & String accuracy on true notes & Fingering alone \\
\bottomrule
\end{tabularx}
\end{table}

\begin{caution}[: comparing numbers]
\begin{itemize}
\item An F50 of 0.80 and a final-tab F1 of 0.80 are not the same thing: one is
  on raw notes with a 50\ms{} window, the other on the finished tab with an
  80\ms{} window.
\item Mean and median differ when a few items fail badly; Python benchmarks
  report means, Rust benchmarks medians.
\item Labels carry timing offsets. Guitar-TECHS labels (from a MIDI pickup) are
  25--40\ms{} late, IDMT labels about 30\ms{} early, GOAT up to $\pm40\ms$
  \measured{} (\code{align_labels.py}). With a 50\ms{} tolerance this alone can
  move F1 a lot, so our benchmarks shift each recording's labels by its median
  offset first.
\item Many published results are on the dataset the model was trained on
  (different files, same players or studio). The use case is zero-shot.
\item Agreement with one human tab is a lower bound on quality: several
  fingerings of the same passage can be equally good.
\end{itemize}
\end{caution}
